\section{Normalizing an incidence flag}
\label{sec:fibers}
An incidence geometry consists of a set of objects, a finite type set $I$, and a reflexive symmetric incidence relation for which distinct objects of the same type are not incident. A flag is a set of pairwise incident objects. We require every flag to extend to a chamber, that is, a flag containing one object of every type. A panel omits exactly one type. The geometry is thin if each panel is contained in exactly two chambers. The residue of a flag consists of the objects incident with it, of types not present in the flag, with the induced incidence relation. Residual connectedness means that every residue of rank at least two, including the whole geometry, has a connected incidence graph.

Let $G$ be a group and let $H_i\leq G$ for $i\in I$. We use left cosets $gH_i$ as objects of type $i$, with incidence given by nonempty intersection. This construction always gives an incidence system; when it is called a geometry below, the flag-extension condition is an explicit hypothesis. The action is left multiplication. Write
\[
 H_J=\bigcap_{j\in J}H_j,\qquad H_\varnothing=G,
 \qquad S_{J,k}=\bigcap_{j\in J}H_jH_k\quad(k\notin J).
\]
The product $H_jH_k$ is a set product; it need not be a subgroup.

\begin{theorem}[The extension fiber]\label{thm:fiber}
Suppose $G$ is transitive on flags of type $J$. Then its orbits on flags of type $J\cup\{k\}$ are naturally in bijection with
\[
                 H_J\backslash S_{J,k}/H_k.
\]
Here the expression denotes the orbits of $H_J$ on the set of cosets $xH_k$ contained in $S_{J,k}$. The stabilizer of the extension represented by $xH_k$ is $H_J\cap xH_kx^{-1}$.
\end{theorem}
\begin{proof}
Move the $J$-part of a flag to $\{H_j:j\in J\}$. A type-$k$ object $xH_k$ meets $H_j$ exactly when $x\in H_jH_k$. Consequently the possible remaining objects are precisely $S_{J,k}/H_k$. Two normalizations of the same $J$-flag differ by an element of its stabilizer $H_J$, so the indicated orbit space is independent of the normalization. Conversely an element relating two normalized extensions fixes their common $J$-part and therefore lies in $H_J$. This proves both directions of the orbit correspondence. Fixing the last object adds exactly the condition $g\in xH_kx^{-1}$ to $g\in H_J$, giving the stabilizer formula.
\end{proof}

\begin{corollary}\label{cor:twofiber}
If $H_J=1$, every normalized extension has trivial stabilizer, and the number of flag orbits is the number of cosets in $S_{J,k}/H_k$. In particular there are exactly two such orbits precisely when
\[
                 S_{J,k}=H_k\ \dot\cup\ xH_k
                 \quad\text{for some }x\notin H_k.
\]
If $H_k$ is finite, this is equivalent to $|S_{J,k}|=2|H_k|$.
\end{corollary}
\begin{proof}
The action of the trivial stabilizer makes each normalized extension an orbit. The identity belongs to every set product, so $H_k$ is one of the cosets. All cosets have size $|H_k|$; when this is finite their number is determined by the displayed cardinality. No such division of infinite cardinals is used.
\end{proof}

The normalization argument is the usual orbit--stabilizer argument applied to an extension of a flag. It contains the two-orbit double-coset calculation in Zhang--Leemans, Theorem 3.2. We use the form with arbitrary stabilizer because it states exactly which hypothesis permits the simpler two-coset test.
